\section{Additional Miscellaneous Results and Discussions}
In this section we list additional insights and results which were found while studying the LP's polytope, some analytical, some via computer analysis.

\subsection{Refining the LP to More Accurately Describe Search Trees}
\label{section_refining_the_LP}
In the discussion leading to Definition~\ref{definition_LP_program}, we noted that the LP does not fully capture the behavior of STTs. In the following, we list a few ``strong'' properties of STTs which were not taken into account, show how to weakly account for them, and finally explain why adding them to the definition of the LP does not remove all the non-integer vertices.

\begin{enumerate}
    \item Path monotonicity: Consider three nodes on a path, in order $a,b,c$. Then $X_{ac} = 1$ only if $a$ is the LCA of itself and $c$, and therefore if $X_{ac} = 1$ then   $X_{ab} = 1$. We get that $X_{ab} \ge X_{ac}$. More generally, if $a,b_1,b_2,\ldots,b_k$ is a path in the topology, $X_{ab_1} \ge X_{ab_2} \ge \ldots \ge X_{ab_k}$. As an example, on the path topology, the ancestry matrix $X$ is monotone-weakly-decreasing in each row, from the main diagonal to the left and to the right. Note that for pure STTs this path monotonicity implies a step function from $0$ to $1$ to $0$ where up to two of these segments may be empty, in particular the row may be constant. In general feasible solutions the steps may be more gradual, but the up-then-down monotonicity is preserved.

    \item Ancestry transitivity: in STTs, if $i$ is an ancestor of $j$ and $j$ is an ancestor of $k$, then $i$ is an ancestor of $k$. This can be written in the non-linear and non-convex form: $X_{ik} \ge \min(X_{ij},X_{jk})$. A linear compromise would be $X_{ik} \ge X_{ij} + X_{jk} - 1$. (While this constraint holds for any triplet, if $k \in (i \leftrightsquigarrow j)$ it is implied by path monotonicity, because then $X_{ik} \ge X_{ij}$, and $X_{jk} \le 1$ for a minimization objective.)

    \item LCA separation: For $k \in (i \leftrightsquigarrow j)$, if $k$ is an ancestor of $i$, then $i$ cannot be an ancestor of $j$. Assume that $i$ is an ancestor of $j$ too, then by the transitivity of ancestry $k$ would be an ancestor of $j$ as well, but since $k$ is between $i$ and $j$, it separates them in contradiction to $i$ being an ancestor of $j$. We can represent this in a linear form as $X_{ki} + X_{ij} \le 1$. Furthermore, we can also require $X_{ji} + X_{ij} \le 1$ for any STT, thus overall: $\forall k \in (i \leftrightsquigarrow j]: X_{ki} + X_{ij} \le 1$.

    \item Refining $Z$: we originally defined $Z_{kij} \le X_{ki},X_{kj}$ as a linearization of $Z_{kij} = \min \{ X_{ki},X_{kj} \} $. Considering STTs for which $X$ has $0/1$ values, we can restrict further and also demand: $Z_{kij} \ge X_{ki} + X_{kj} - 1$.

    \item Row-Min/Column-Max: Consider the ancestry matrix $X$ of some STT with root $r$. The minimum in each row $i \ne r$ is $0$ because $X_{ir} = 0$, and the minimum in row $r$ is $1$ because $\forall j: X_{rj} = 1$. Similarly, the maximum in each column $i \ne r$ is $1$, and $0$ in column $r$. We can define new variables $m_i$ for the minimum in row $i$ and $M_i$ for the maximum in column $i$, and add the constraints: $\forall j: X_{ji} \le M_i$, $\forall j: m_i \le X_{ij}$ and $\sum_{i}{M_i} = n-1$, $\sum_{i}{m_i} = 1$. (The counterparts max-row and min-column are less useful because they change for $i \ne r$ depending on the STT.)
\end{enumerate}

So far it looks very promising, few natural constraints of STTs were missing from the LP, which we can add. Moreover, we introduced only $O(|X|) + |Z|$ extra inequalities, so the problem is still polynomial.\footnote{$|X|$ and $|Z|$ denote the numbers of $X$ and of $Z$ variables, respectively.} However, these extra inequalities do not fix the polytope, and non-integer vertices remain plentiful where they exist. It may remove some, and create others (due to new hyperplanes), but it does not remove all of them which is what we hoped for, if this was an extended formulation. To see this, we only need to go back to our main example, the vertex $P$ in Figure~\ref{figure_explicit_counter_example}. It satisfies all the additional constraints, as we prove:

\begin{enumerate}
    \item Path monotonicity: $X_{12} = X_{13} = \frac{1}{2} > X_{14} = X_{15} = X_{16} = X_{17} = 0$; $X_{23} = 1 > X_{24} = X_{25} = X_{26} = \frac{1}{2} > X_{27} = 0$; $X_{3i} = 0$ for all $i \ne 3$; $X_{43} = 1 > X_{42} = X_{41} = X_{46} = \frac{1}{2} > X_{47} = 0$; $X_{54} = X_{53} = \frac{1}{2} > X_{52} = X_{51} = X_{56} = X_{57} = 0$; $X_{63} = 1 > X_{6i} = \frac{1}{2}$ for all $i \ne \{3,6\}$; $X_{7i} = \frac{1}{2}$ for all $i \ne 7$.
    
    \item Ancestry transitivity: note that if $X_{ij},X_{jk} \le \frac{1}{2}$ then the inequality becomes trivial because $X_{ik} \ge 0$ anyway. So the only entries that can possible make $P$ infeasible are pairs that include at least one of the entries $X_{23}=X_{43}=X_{63}=1$. If $i \in \{2,4,6\}$ and $j=3$ then $X_{jk}=0$ regardless of $k$ and we get the constraint $X_{ik} \ge 0$ (always true). If $j \in \{2,4,6\}$ and $k=3$ we need to check that $X_{i3} \ge X_{ij}$ for any $i \notin \{j,3\}$. This holds because $X_{i3} \ge \frac{1}{2} \ge X_{ij}$ (here $j \ne 3$).
        
    \item LCA separation: If $X_{ki},X_{ij} \le \frac{1}{2}$ then $X_{ki} + X_{ij} \le 1$ is satisfied. Then, to verify that $P$ is feasible, we only need to consider the case that $i=3$ or $j=3$. If $i=3$ then $X_{ij} = 0$ for any $j$ and the inequality is satisfied. For $j=3$ we consider $i \in \{2,4,6\}$, which is a neighbor of $j$, hence $k = j = 3$ is the only option for $k$, and the inequality holds.

    \item Refining $Z$: If $X_{ki},X_{kj} \le \frac{1}{2}$ then $Z_{kij} \ge 0$ is trivial. Then focus on $i=3$ ($j=3$ is symmetric), with $k \in \{2,4,6\}$ and the inequality reduces to $Z_{k3j} \ge X_{kj}$. The requirement $k \in (i \leftrightsquigarrow j)$ together with $i=3$ implies that $(k,j) \in \{ (2,1), (4,5), (6,7) \}$. Then we can verify that $Z_{213}  = X_{21} = Z_{435} = X_{45} =Z_{637}= X_{67}   = \frac{1}{2}$ hence the inequalities hold.

    \item Row-Min/Column-Max: the min-row vector is $\vec{q} = [0,0,0,0,0,\frac{1}{2},\frac{1}{2}]$, hence we can choose $\vec{m} = \vec{q}$ to satisfy the row-min part. The max-column vector is $\vec{Q} = [\frac{1}{2},\frac{1}{2},1,\frac{1}{2},\frac{1}{2},\frac{1}{2},\frac{1}{2}]$, and there are plenty of choices for $M$ such that $\forall i: M_i \ge Q_i$ and $\sum_i{M_i} = n-1$, which satisfy the column-max part.
\end{enumerate}

In conclusion, even though we refined the LP to capture more accurately the behavior of STTs, we did not get an extended formulation that removes every non-STT vertices, and an integrality gap remains. One reason may be that we did not capture all possible refinements. Indeed, when we compute the convex hull of the $X$s induced by STTs, there are many additional constraints. To name just one such family, for every permutation $\pi$ without fixed points ($\forall i: \pi_i \ne i$), $\sum_i{X_{i,\pi_i}} \ge 1$. The intuitive interpretation is that some node must be the root $r$ of the STT, therefore $X_{r,\pi_r} = 1$. The size of this set is roughly $\frac{n!}{e}$.






\subsection{Eliminating the \emph{Z} Variables from the Formulation}
\label{subsection_with_and_without_z_fourier_motzkin}
In the LP definition in Section~\ref{subsection_linear_program}, we used $Z_{kij}$ variables whose conceptual purpose is to represent an LCA relationship of $k$ over $i$ and $j$, practically a linear replacement to a term $\min\{X_{ki},X_{kj}\}$ which cannot simply be used in an LP formulation. It is natural to ask whether we can eliminate these variables and replace them by additional constraints on the $X_{ij}$ variables. This may have drastic effects, in particular making many integer vertices of the original LP polytope non-vertices, as discussed in Remark~\ref{remark_int_vertices_disappear_when_no_Z}. To get rid of all the $Z$ variables, we can take every constraint that  contains $d$ minimum-terms (or rather $d$ $Z$-variables each representing the minimum of a pair of $X$-variables), and replace it by $2^d$ inequalities where each inequality corresponds to a choice of a particular $X$-variable from each  minimum-pairs as in Example~\ref{example_Z_elimination}.\footnote{We can formally and compactly define all the inequalities as follows: Let $G$ be the set of $2^d$ functions from $(i \leftrightsquigarrow j)$ to $\{i,j\}$. Then for each $g \in G$ we get the inequality $X_{ij}+X_{ji} + \sum_{k \in (i \leftrightsquigarrow j)}{X_{k,g(k)}} \ge 1$.}
Even though we get an exponential number of constraints, we can still use a separation oracle (cutting-plane oracle) to solve the LP in polynomial time, since identifying the tightest inequality out of each such group of $2^d$ inequalities that replace a single original inequality is easy:
 just take the inequality where the smallest of each pair $X_{ki}$ and $X_{kj}$ participates.

\begin{example}
\label{example_Z_elimination}
Consider an ancestry constraint on a pair $a$ and $d$ with two nodes $b,c$ on the path between them: $X_{ad} + X_{da} + \min\{ X_{ba},X_{bd} \} + \min\{ X_{ca},X_{cd} \} \ge 1$. The constraint is satisfied if and only if the following set of four constraints is satisfied:
$X_{ad} + X_{da} + X_{ba} + X_{ca} \ge 1$ and
$X_{ad} + X_{da} + X_{ba} + X_{cd} \ge 1$ and
$X_{ad} + X_{da} + X_{bd} + X_{ca} \ge 1$ and
$X_{ad} + X_{da} + X_{bd} + X_{cd} \ge 1$.
\end{example}

The Fourier-Motzkin elimination technique~\cite{BookForFourierMotzkin} can also be used to eliminate the $Z$ variables. One can verify that its application yields the same set of inequalities.

Observe that once we eliminate the $Z$ variables, ancestry constraints must be inequalities. Previously, we could make them equalities $(=1)$, but now we cannot do this since we do not know which among the $2^d$ constraints is tight. Overall, the trade-off of eliminating $Z$ gives us a more complex presentation of the polytope (more inequalities) in exchange for less dimensions (no $Z$) as well as less vertices overall as it turns out.

\begin{remark}
\label{remark_int_vertices_disappear_when_no_Z}
It is interesting to mention that the proof of Theorem~\ref{theorem_non_tree_vertex} breaks under the new formulation of the LP without $Z$ variables. There, the second case in the proof completely disappears because there are no $Z$ variables to manipulate. The first case of cyclic ancestry is feasible, but it might not be a vertex anymore. For an example of vertices that become non-vertices without $Z$, revisit Table~\ref{table_vertices_n3}: The two ``LCA abuse'' vertices become dominated by the vertex that corresponds to the STT whose root is $2$ when we eliminate $Z$. More generally, the set $\mathbb{P}'$ as defined in the proof of Theorem~\ref{theorem_non_tree_vertex} gets much smaller or even empty, and no longer contains all the STT induced vertices.
\end{remark}



%%% A DETAILED DISCUSSION ON BLOWUP FACTOR, FOR COMPLETENESS.
%Since $|Z| = O(n^3)$ (more accurately, $|Z| \le \binom{n}{3}$) while $|X| = \binom{n}{2}$, it may be worthwhile to eliminate the $Z$ variables to reduce the dimension of the LP. On the other hand, in exchange, we increase the number of inequalities by a factor of $O(2^d)$ where $d$ is the diameter of the topology. Note that if $d = O(1)$ then we only have $Z = O(n^2)$ to begin with. Overall, the trade-off of eliminating $Z$ gives us a more complex presentation of the polytope (more inequalities) in exchange for less dimensions (no $Z$) as well as less vertices overall as it turns out (see also Remark~\ref{remark_int_vertices_disappear_when_no_Z}). For an explicit example how vertices become non-vertices in the new formulation, revisit Table~\ref{table_vertices_n3} of all the vertices for the topology with $3$ nodes. The two ``LCA abuse'' vertices become dominated by the vertex that corresponds to the STT whose root is $2$ when we eliminate $Z$. As a quick example of the blow-up factor, consider the two extreme cases: In a path topology over $n$ nodes, the number of ancestry constraints grows by a factor of $\approx 2^{n+1}/n^2$ (the long paths blow-up exponentially, but most paths are short), and in a star topology the blow-up factor is only $2(1-1/n)$.

%%% THE EXPLICIT ANALYSIS FOR THE BLOWUP FACTOR FOR PATHS AND STARS, FOR COMPLETENESS.
% \begin{example}
% \label{example_elimination_eq_blowup}    
% Eliminating $Z$ for a path topology over $n$ nodes, blows-up the number of ancestry-inequalities by a factor of $\approx 2^{n+1}/n^2$. Indeed, an inequality with $d$ eliminations results in $2^d$ inequalities. We have $n-h$ pairs at edge-distance $h$ from each other. Therefore, the averaged blow-up factor is:
% $$\frac{\sum_{h=1}^{n-1}{(n-h) \cdot 2^{h-1}}}{\binom{n}{2}} = \frac{2^n-n-1}{n(n-1)/2} \approx \frac{2 \cdot 2^n}{n^2}$$

% Eliminating $Z$ for a star topology over $n$ nodes, blows-up the number of ancestry-inequalities by a factor of $2 \big( 1 - \frac{1}{n} \big)$. Indeed, in $n-1$ pairs where the center-node is part of a pair, we do not have any $Z$. The rests $\binom{n-1}{2}$ pairs blow-up by a factor of $2$ due to a single $Z$ variable in each. Therefore, the averaged blow-up factor is:
% $$\frac{(n-1) + 2 \binom{n-1}{2}}{\binom{n}{2}} =
% \frac{1 + (n-2)}{\frac{n}{2}} =
% 2 \Big( 1 - \frac{1}{n} \Big)
% $$
% \end{example}

Overall, it is probably better to eliminate the $Z$ variables from the LP. It is ``more true'' (by capturing the minimum as intended), and it also affect on the density of non-integer vertices, as we explain next. As a warm-up, note that given that some nodes disappear, the relative density of integer versus non-integer nodes may change. We can refine this differentiation by categorizing each vertex according to its highest denominator. In this case, consider for example the path topology over $n=5$ nodes. By enumerating the vertices of its corresponding polytope, we find that without eliminating $Z$, there are $5983$ integer vertices (out of which only $42$ are due to STTs), $3886$, half-integer vertices, and $76$ third-integer vertices. Compactly, we write $[5983,3886,76]$.\footnote{There are no vertices with both half-integer and third-integer coordinates. This scenario is possible in vertices for larger topologies, by Theorem~\ref{theorem_combine_trees}.
% see Example~\ref{example_mixed_halves_and_thirds}.
Also, non of these non-integer vertices is projected onto a vertex in $D$-space, so in $D$-space they are all non-vertices.} Counting the vertices of the polytope in the $Z$-eliminated version, we get much less vertices, distributed $[519,158,7]$ (still $42$ STT vertices, out of the $519$). Observe that, for example, the percentage of non-integer vertices drops from $40\%$ to $24\%$, and that third-integer vertices rise from $0.76\%$ to $1.02\%$. Changed percentages may affect our probability of finding vertices
when we use a random direction as the objective.
%\footnote{A uniformly random direction may not result in a uniformly random vertex, but the discussion is mostly informal at this point anyway.}

\begin{table}[!ht]%[!ht]
    \scriptsize
    \begin{center}
    \begin{tabular}{|c|l|l||l|l||c|c|}

    \hline
    \multicolumn{1}{|c|}{} & \multicolumn{2}{|c||}{$(X,Z,D)$-Direction} & \multicolumn{2}{|c||}{$(X,D)$-Direction} & \multicolumn{2}{|c|}{$D$-Direction} \\
    \hline
    Path length ($n$) &
    LP &
    $Z$-elim. &
    LP &
    $Z$-elim. &
    LP &
    $Z$-elim. \\
    \hline
    \hline
    6 & [1, 2] & [1, 2] & [1, 2] & [1, 2] & [1] & [1] \\
    \hline
    7 & [1, 2, 3, 4] & [1, 2] & [1, 2] & [1, 2] & [1] & [1] \\
    \hline
    8 & [1, 2, 3, 4, 5] & [1, 2, 3] & [1, 2, 3] & [1, 2, 3] & [1] & [1] \\
    \hline
    9 & [1, 2, 3, 4, 5] & [1, 2, 3] & [1, 2, 3] & [1, 2, 3] & [1] & [1] \\
    \hline
    10 & [1, 2, 3, 4, 5, 6] & [1, 2] & [1, 2] & [1, 2] & [1] & [1] \\
    \hline
    11 & [1, 2, 3, 4, 7] & [1, 2, 3] & [1, 2, 3] & [1, 2, 3] & [1] & [1] \\
    \hline
    12 & [1, 2, 3, 4, 5, 6, 7, 8, 9, 11, 12, 14] & [1, 2] & [1, 2] & [1, 2] & [1] & [1] \\
    \hline
    13 & [1, 2, 3, 4, 5, 6, 7, 8, 9, 11, 12, 16, 20] & [1, 2] & [1, 2] & [1, 2] & [1] & [1] \\
    \hline
    14 &  [1, 2, 3, 4, 5, 6, 8, 9, 10,11, 12, 13,18, 20, 23, 43] & [1, 2] & [1, 2] & [1, 2] & [1] & [1] \\
    \hline
    15 & \begin{tabular}{@{}l@{}}[1, 2, 3, 4, 5, 6, 7, 8, 9, 10, 11, 12, 14, 16, 20, \\ 22, 23, 28, 40, 47, 80] \end{tabular} & [1, 2, 3] & [1, 2, 3] & [1, 2, 3] & [1] & [1] \\
    \hline
    16 & 
    \begin{tabular}{@{}l@{}l@{}}[1, 2, 3, 4, 5, 6, 7, 8, 9, 10, 11, 12, 13, 14, 15, \\ 17, 18, 20, 22, 23, 24, 26, 28, 34, 36, 39, 46, 58, \\ 67, 87, 92, 111, 134, 174, 257, 261, 522, 999] \end{tabular}
     & [1, 2, 3] & [1, 2, 3] & [1, 2, 3] & [1] & [1] \\
    \hline
    \end{tabular}
    \end{center}
    \caption{\footnotesize{A summary of denominators of vertex coordinates found when solving the LP for path topologies, of different lengths, in six variations. The vertices were found by sampling directions, which is not exhaustive as can be seen since for $n=5$ we know of third-integer vertices in $(X,Z,D)$-space, but here for $n=6$ we have not found such vertices. We chose the objective (weights) in either: (1) a general direction with arbitrary weights for $(X,Z,D)$, (2) $(X,D)$-direction with weights $0$ for all $Z$, and (3) depths-only direction with non-zero weights only for $D$. In each of these three methods, we solve the LP as in Definition~\ref{definition_LP_program} (LP column), and its variant where we eliminate $Z$ as in Section~\ref{subsection_with_and_without_z_fourier_motzkin} ($Z$-elim. column). It is clear that an arbitrary direction reveals fractions with increasing denominators, yet they seem to be artifacts of the presence of the $Z$, and are either non-existent or very rare when eliminating $Z$ or when considering $(X,D)$ or $D$ directions. Furthermore, all the vertices found in $D$-directions are integer. Among the last five columns, $2$ and $3$ denominators never appear together, i.e., every vertex is either integer, half-integer or third-integer. However, this fact is an artifact of sampling the rare mixed cases, which exist by Theorem~\ref{theorem_combine_trees} (e.g., combine two paths of length $n=5$ with an extra node to get a path of $n=11$).}}
    \label{table_path_tree_fractions}
\end{table}

We cannot enumerate all the vertices in larger topologies, and the interesting question is whether the ``\emph{fractionality}'' of the polytope is affected only by changed percentages, or whether some types of fractions completely disappear or being introduced when we eliminate $Z$. By solving the LP in random directions, it looks like there is indeed a strong difference in behavior. We focused our numerical-analysis on topologies that are paths, see Table~\ref{table_path_tree_fractions}. There, when solving the LP with objectives in arbitrary directions ($(X,Z,D)$ may all have positive weights) we find that the denominators grow when the topology grows. However, when the objective is only in $(X,D)$ directions, 
which are the only direction that remain when
we eliminate $Z$ from the LP, the fractionality only contains integer, half-integer and third-integer vertices. Finally, when solving either version of the LP in a pure $D$-direction, which is the ``intended'' kind of objective, we only find integer vertices. These results hint that something is inherently ``less fractional'' when we eliminate $Z$, even if they are based on sampling rather than a full enumeration or an analytic proof. These results also hint at the possibility that the LP perhaps finds an optimal STT when path topologies are considered.


% %%% Omitted for the sake of paper's length. Keep it in the comments, for the curious readers who want examples of the non-integer half-integer and third-integer vertices.
% \begin{example}
% \label{example_non_int_vertex_dominated_in_D_space}
% Consider the path with $n=5$ vertices. As mentioned in Section~\ref{subsection_with_and_without_z_fourier_motzkin} it has a total of $165$ non-integer vertices when we eliminate $Z$ from the LP. Equation~(\ref{equation_path5_vertices}) shows two such vertices, a half-integer and a third-integer, such that both are non-vertices when projected into $D$-space. We present them in a form similar to Figure~\ref{figure_explicit_counter_example}, $X$ as a matrix, and $D$ below it (column sums).

% {\scriptsize
% \begin{equation}
% \label{equation_path5_vertices}
% % vertex #165 among the non-integer vertices when enumerated by the script.
% P_1 = \frac{1}{2} \cdot 
% \begin{bmatrix}
% . & 1 & 0 & 0 & 0 \\
% 1 & . & 1 & 1 & 0 \\
% 1 & 1 & . & 0 & 1 \\
% 1 & 1 & 2 & . & 1 \\
% 0 & 0 & 0 & 1 & . \\
% \hline
% 3 & 3 & 3 & 2 & 2 \\
% \end{bmatrix}
% ~
% % vertex #162 among the non-integer vertices when enumerated by the script.
% P_2 = \frac{1}{3} \cdot 
% \begin{bmatrix}
% . & 1 & 0 & 0 & 0 \\
% 2 & . & 2 & 3 & 1 \\
% 1 & 2 & . & 0 & 2 \\
% 1 & 0 & 3 & . & 1 \\
% 0 & 0 & 0 & 2 & . \\
% \hline
% 4 & 3 & 5 & 5 & 4 \\
% \end{bmatrix}    
% \end{equation}
% }

% $P_1$ is not a vertex in $D$-space, since it is strictly worse than the convex combination of STT depths $\frac{1}{2}([2,1,2,0,1] + [1,2,0,2,1])$. $P_2$ is not a vertex in $D$-space, since it is strictly worse than the convex combination of STT depths $A \equiv \frac{1}{3}([1,0,2,1,2]+[2,1,2,0,1]+[1,2,0,2,1])$, and we can represent it as a convex combination of $A$ and positive rays.

% Interestingly, all $165$ non-integer vertices are non-vertices when projected to $D$-space, non of them coincides with an integer vertex, and non of them even has all-integer depths. To keep this example short, we leave out the verification that these points are feasible and are vertices, and just claim to have found them via coded enumeration (Section~\ref{section_code} discusses the code).
% \end{example}





% %%% The following subsection is a bit ``divergent''. It essentialy discusses, in a slightly verbose manner, the attempt to enforce that we do not get third-integer vertices where ancestry inequalities overflow the sum of 1. In summary it is hard to enforce, and a bit artificial. I kept the text commented-out.
% \subsection{Tightening Third-integer Vertices?}
% \label{subsection_coordinates_fractions}
% In Section~\ref{subsection_with_and_without_z_fourier_motzkin} we discussed the fact that even when we eliminate that $Z$ variables, which likely add vertices with values that are not ``nice'', half-integer and third-integer vertices remain.

% A deeper look into the third-integer vertices of the path graph with $n=5$ nodes reveals that there is at least one pair of nodes $i$ and $j$ such that $X_{ij} = X_{ji} = \frac{2}{3}$, see Example~\ref{example_non_int_vertex_dominated_in_D_space}. This is feasible because we weakened the ancestry constraint to be an inequality. One may consider to restrict further and demand that $X_{ij} + X_{ji} \le 1$ for every pair (holds for STTs). Does it remove the third-integer vertices? In general, the answer is negative. These extra constraints do remove all third-integer vertices for the path of length $n=5$. In return, the number of vertices in total, integer and half-integer, increases. However, when testing random objective directions for larger paths, we find that already for $n=8$ that satisfy the extra constraint (possibly also for $n=6,7$). We provide one such vertex in Figure~\ref{figure_third_int_vertex_path8} (without a proof of it being a vertex).

% This subsection might be a moot point since initially the ancestry constrains were equalities which we relaxed, and $X_{ij} + X_{ji} \le 1$ does not fully turns us back to an equality. This ``overflow'' can be seen in Figure~\ref{figure_third_int_vertex_path8} for example: $X_{36}+X_{63}+\min(X_{43},X_{46})+\min(X_{53},X_{56}) = 0+\frac{2}{3}+\frac{2}{3}+0 > 1$. The issue with trying to stick to ancestry equalities is that then we cannot eliminate the $Z$ variables (when we do, we get back to inequalities). 


% \begin{figure}[ht]
%     {\scriptsize
%     \begin{equation}
%     \label{equation_third_int_vertex_path8-6_vertices}
%     % vertex #165 among the non-integer vertices when enumerated by the script.
%     P = \frac{1}{3} \cdot 
%     \begin{bmatrix}
%     . & 2 & 2 & 0 & 0 & 0 & 0 & 0 & \\
%     1 & . & 2 & 2 & 1 & 0 & 0 & 0 & \\
%     0 & 1 & . & 0 & 1 & 0 & 0 & 0 & \\
%     2 & 1 & 3 & . & 2 & 2 & 1 & 1 & \\
%     0 & 0 & 0 & 1 & . & 1 & 0 & 0 & \\
%     1 & 2 & 2 & 0 & 2 & . & 2 & 2 & \\
%     1 & 0 & 0 & 2 & 1 & 1 & . & 3 & \\
%     0 & 0 & 0 & 0 & 0 & 0 & 0 & . & \\
%     \hline
%     5 & 6 & 9 & 5 & 7 & 4 & 3 & 6 \\
%     \end{bmatrix}
%     ~
%     % vertex #162 among the non-integer vertices when enumerated by the script.
%     weights = 
%     \begin{bmatrix}
%     . & 0 & 2 & 11 & 2 & 21 & 6 & 22 & \\
%     22 & . & 21 & 6 & 2 & 14 & 20 & 4 & \\
%     22 & 17 & . & 24 & 11 & 20 & 16 & 11 & \\
%     1 & 14 & 7 & . & 17 & 7 & 17 & 24 & \\
%     18 & 17 & 21 & 15 & . & 0 & 3 & 21 & \\
%     8 & 10 & 0 & 19 & 2 & . & 11 & 17 & \\
%     4 & 16 & 17 & 2 & 8 & 16 & . & 1 & \\
%     14 & 1 & 15 & 18 & 7 & 5 & 1 & . & \\
%     \hline
%     18 & 14 & 9 & 9 & 12 & 13 & 17 & 9 \\
%     \end{bmatrix}    
%     \end{equation}
%     }
    
%     \caption{\footnotesize{A vertex $P = (X,D)$ of the LP for the path of $n=8$ nodes, that is third-integer, such that $\forall i,j: X_{ij} + X_{ji} \le 1$. $weights$ is the objective direction for which the vertex can be found. The location of each weight corresponds to the variable's location, note that this is an $(X,D)$-direction.}}
% \label{figure_third_int_vertex_path8}
% \end{figure}


\subsection{Are Path-monotonicity Constraints Helpful?}
In Section~\ref{section_refining_the_LP} we listed several possible refinements of the LP. Among them, the set of path-monotonicity constraints is arguably one of the most natural and simple to grasp. We have already seen that adding these constraints does not remove all non-integer vertices, but one can wonder whether it helps. It turns out to be a double-edged sword.

At a first glance, it has benefits. The small topologies with non-integer vertices that we found have $74$ such vertices in total after sieving symmetric copies (revisit Table~\ref{table_trees_summary_analysis_vertices_etc_FEW}). For each specific topology, more than half satisfy these constraints, so no topology becomes ``truly nicer''. Testing singular vertices is sort of cheating, because we need to see how the polytope changes as a whole. When considering the path over $5$ nodes, we get promising results. Consider the version without $Z$ variables, as discussed in details in Section~\ref{subsection_with_and_without_z_fourier_motzkin}. When we add path-monotonicity constraints, we still have non-integer vertices overall, but instead of having $165$ as such ($158$ half-integer and $7$ third-integer), the new polytope only has $4$ non-integer vertices, all half-integer. This is not quite fully-integer, but it looks like a promising simplification.

Alas, some topologies are badly affected. Consider topology $U_{(5,1)}$ ($5$ nodes, ``T''-shaped) which was found to have an integer polytope for the original LP. If we add path-monotonicity constraints, its $(X,Z,D)$-polytope now has half-integer and third-integer vertices. So, to conclude this discussion: while path-monotonicity seem like a clean and simple addition to the LP, it is not clear whether it helps or complicates things.


\subsection{Multiple Vertices with the Same Depths}
Throughout the paper, we sometime discuss vertices in $(X,Z,D)$-space, and sometimes in $D$-space. In particular, the normals method uses $D$-space to solve and find ``$(X,Z,D)$-vertices'', which we then project back to $D$-space. It is natural to ask whether the projection to $D$-space results in collisions of vertices, that is, if there exists $(X,Z,D)$-vertices, $P \ne P'$, such that $P_D = P'_D$. Because the LP without $Z$ variables is arguably nicer and ``more true'' (recall Section~\ref{subsection_with_and_without_z_fourier_motzkin}), let us wonder about collisions of vertices of the $(X,D)$-polytope of the LP when projected to $D$-space.

For STTs the answer is negative: given a depths-vector $D$ induced by an STT, there is exactly one $(X,D)$-vertex in its pre-image. Indeed, the depths-vector $D$ uniquely implies the value of all the $X$ variables, because there is a unique index $r$ (the root) such that $D_r=0$, so it must be that $\forall i \ne r: X_{ir}=0,X_{ri}=1$. Then we recurse in each connected component, now looking for a vertex of depth one in each connected component etc.

In contrast, general vertices may collide. When we solved the LP without the $Z$ variables, in all the primary directions that discover non-STT vertices, we found a few collisions. Recall that the objectives that we used had non-zero coefficients only for the $D$-variables, so to truly study the collisions one should use general objectives. Furthermore, note that refinements of the LP such as those discussed in Section~\ref{section_refining_the_LP} may affect these collisions.
% [Slightly more information:] The collisions that we find are only of pairs, but in general there might be collisions of higher cardinality. The collisions that we find (depends on the solver being used since the objective function wsa in a pure $D$-direction): Two collisions for the topology $U_{(8,5)}$, one for $U_{(8,6)}$, and eight for $U_{(8,13)}$. Note that these findings only prove that collisions happen and say nothing about other collisions that we missed or the cardinality of these collisions (may be more than just pairs).

